% Chapter 1: Introduction
% Contains: Overview, Target Domains, Ecosystem Architecture, System Architecture

% ============================================================
\section{Overview}
% ============================================================

Avena is a distributed fleet management system for vehicle-centric edge computing environments. It manages devices---tractors, irrigation controllers, field sensors, traffic signal controllers, connected vehicles---operating where network connectivity is intermittent, expensive, or topologically unstable.

\textbf{Core thesis:} Existing pub/sub systems (NATS, Kafka, MQTT) assume reliable networks. DTN solutions assume point-to-point messaging. Avena bridges both: subject-based pub/sub semantics with DTN-style store-and-forward delivery and cost-aware routing.

Applications use a NATS-compatible pub/sub API. The infrastructure handles intermittent connectivity, mobile relays, heterogeneous link costs, and multi-path routing transparently. Code running on well-connected cloud nodes and intermittently-connected field sensors uses the same API.

\subsection{Design Principles}

\begin{enumerate}[label=\textbf{\arabic*.}]
    \item \textbf{Disconnection is normal.} Devices operate autonomously when disconnected and synchronize state opportunistically when connectivity exists.
    \item \textbf{Mobility creates topology.} Mobile assets (tractors, trucks, vehicles) serve as data carriers between disconnected network segments, forming temporal paths through the network.
    \item \textbf{Security without central authority.} A delegated certificate model enables authorization decisions without requiring connectivity to a central server.
    \item \textbf{Transparent edge-to-cloud.} Application code is identical on edge devices and cloud servers; the infrastructure abstracts network complexity.
\end{enumerate}

% ============================================================
\section{Target Domains}
% ============================================================

Avena addresses two primary application domains that share vehicle-centric networking challenges.

\subsection{Agricultural Edge Computing}

Modern farms deploy diverse computational assets across thousands of acres: tractors with guidance systems and telematics, irrigation controllers responding to soil moisture sensors, autonomous field robots, stationary weather stations, and grain handling equipment with embedded controllers. Network infrastructure is sparse---a sensor deployed in a remote field section may have no persistent connectivity. Cellular coverage, where available, is expensive and unreliable.

The operational pattern is vehicle-centric: tractors, combines, and grain carts traverse the operation, creating transient connectivity with fixed sensors and with each other. A combine harvesting a remote field accumulates yield data with no backhaul; when it returns to the grain elevator, accumulated data synchronizes. A tractor passing a soil moisture sensor downloads readings that have been buffered for days.

Agricultural fleets are typically single-domain (one farm operation), but multi-domain scenarios arise naturally. Custom harvesters bring their own equipment onto client farms. Agronomic consultants need access to field data without full fleet privileges. Neighboring operations may cooperate: a sensor near a property boundary could provide valuable microclimate data to the adjacent farm.

\subsection{Vehicle-to-Everything (V2X) Communication}

Connected vehicle infrastructure presents a complementary challenge profile. Traffic signal controllers, roadside units (RSUs), and vehicles must exchange information for safety applications (collision warnings, signal phase and timing, emergency vehicle preemption) and efficiency applications (traffic flow optimization, platooning, predictive routing). Vehicles traverse the road network at highway speeds, creating rapidly-changing topology.

V2X inherently involves multiple administrative domains with distinct trust requirements:

\begin{description}
    \item[Infrastructure domain:] Traffic signal controllers, roadside units, variable message signs---managed by state or municipal transportation departments.
    \item[Commercial fleet domain:] Delivery vehicles, transit buses, taxis, freight trucks---managed by fleet operators.
    \item[Consumer domain:] Individual vehicles with their own trust boundaries.
\end{description}

V2X motivates several Avena capabilities particularly strongly. Geo-spatial interests are essential: a vehicle cares about signal timing for intersections on its current route, not all signals in the metropolitan area. Priority-aware routing matters because safety messages (Basic Safety Messages, collision warnings) require delivery within 100ms, while less urgent data can tolerate seconds of delay.

\subsubsection{V2X Use Cases}

\textbf{Emergency vehicle preemption.} An ambulance approaching an intersection broadcasts a preemption request. The traffic signal controller, subscribed to emergency vehicle subjects within its geographic vicinity, receives the request and adjusts signal timing. The preemption request is high-priority, low-latency, geographically bounded, and crosses domain boundaries.

\textbf{Cooperative platooning.} A convoy of freight trucks coordinates tight following distances. Each truck subscribes to the position, velocity, acceleration, and braking status of convoy members with high priority and tight latency bounds.

\textbf{Traffic signal optimization.} Traffic signals collect approach data from vehicles. Vehicles approaching a signalized intersection receive Signal Phase and Timing (SPaT) messages. The subscription is geographically bounded (vehicles within 500 feet of the intersection) and time-sensitive.

% ============================================================
\section{Ecosystem Architecture}
% ============================================================

Avena's technical architecture enables a market structure where hardware, software, and operational concerns decouple---creating network effects that benefit all participants.

\subsection{Hardware-Software Decoupling}

Traditional agricultural and V2X systems exhibit tight coupling: hardware vendors build vertically-integrated stacks where equipment requires vendor-specific software, and software assumes specific hardware platforms. This coupling constrains both markets. Hardware value derives solely from the vendor's software investment; software addressable market is limited to that vendor's installed base.

Avena inverts this model. Any workload conforming to the Avena API (NATS-compatible pub/sub with certificate-scoped authorization) deploys to any Avena-compliant device. Hardware vendors compete on physical capabilities---sensing, actuation, power efficiency, environmental hardening---without building complete software stacks. Software vendors compete on application value without manufacturing hardware or negotiating per-OEM integrations.

\textbf{Market dynamics:}
\begin{itemize}[nosep]
    \item \textbf{Hardware vendors} gain value from the software ecosystem they did not build. A sensor manufacturer's device becomes more valuable because third-party analytics, visualization, and control software deploys without integration work.
    \item \textbf{Software vendors} access the combined installed base of all Avena-compliant hardware. An agronomic analytics application deploys to equipment from multiple manufacturers without per-OEM porting.
    \item \textbf{Operators} select hardware and software independently based on operational requirements, avoiding lock-in to either dimension.
\end{itemize}

This parallels the Android ecosystem model: device manufacturers specialize in hardware while application developers target a common platform. The difference is that Avena targets industrial edge computing where workload isolation, offline operation, and delay-tolerant networking are primary concerns rather than user interface consistency.

\textbf{Enabling mechanisms:} The workload model (Section~\ref{sec:workloads}) provides the deployment abstraction. Certificate-based authorization (Section~\ref{sec:certificates}) enables cross-organization trust. The transparent API ensures application code is portable across device classes.

\subsection{Interoperability Through Message Semantics}

Traditional agricultural and fleet systems exchange data through proprietary file formats, vendor-specific APIs, and complex aggregated data structures. Integrating two systems requires either vendor cooperation (unlikely between competitors) or brittle translation layers that break when either system evolves.

Avena shifts the integration model:

\textbf{Atomic messages over aggregated formats.} Rather than exchanging complete data models, systems publish small, self-describing messages on content-addressed subjects. A soil moisture reading is a single message with a timestamp, location, and value---not an entry in a proprietary database export. Schema agreement for atomic messages is tractable; schema agreement for complex nested structures rarely succeeds across organizational boundaries.

\textbf{Consumer-side integration.} Publishers emit messages without assumptions about how consumers will use them. Each consumer integrates message streams into whatever data model serves their use case. The same soil moisture stream feeds a real-time irrigation controller, a seasonal analytics platform, and a regulatory compliance archive---each with its own internal representation.

\textbf{Subject-based discovery.} Content-addressed subjects (e.g., \texttt{app.telemetry.soil.field-north.sensor-7}) provide semantic meaning without requiring schema registries or format negotiation. New consumers discover available data by subscribing to subject patterns; no integration contracts required.

\textbf{Implications for system integration:}
\begin{itemize}[nosep]
    \item Operators integrate systems from non-cooperating vendors by subscribing both to relevant subject patterns and transforming messages in application-specific ways.
    \item Adding a new data consumer requires no changes to existing publishers.
    \item Data portability is inherent: messages are self-describing and flow to any subscriber with appropriate authorization.
    \item Lock-in decreases: switching a component (sensor, analytics platform, visualization tool) requires only redirecting subscriptions, not data migration or format conversion.
\end{itemize}

The messaging architecture (Section~\ref{sec:messaging}) and subject namespace (Section~\ref{sec:subjects}) provide the technical foundation. The ecosystem benefit is that interoperability emerges from architectural choices rather than requiring explicit vendor cooperation.

% ============================================================
\section{System Architecture}
% ============================================================

Avena comprises three interrelated subsystems that together provide delay-tolerant pub/sub messaging and fleet management for vehicle-centric networks.

\textbf{Avena Overlay Network.} The foundational layer provides encrypted, authenticated IP connectivity over heterogeneous physical links. Each Avena device creates a \emph{dedicated Wireguard interface for each (peer, underlay) combination}---if device A can reach device B over both WiFi and cellular, A maintains two separate Wireguard interfaces to B. This enables true multi-path connectivity where Babel routing (RFC 8966) sees each physical path as a distinct link with its own metrics. Babel selects routes based on link quality (latency, loss, bandwidth) and can failover between paths when physical connectivity changes. A parallel low-rate channel (LoRa) provides control-plane connectivity when high-rate links are unavailable, enabling presence announcements and topology updates before IP tunnels establish.

The overlay answers: \emph{can packets flow between devices?} It abstracts physical link heterogeneity, presenting applications with a unified IP network regardless of whether the underlying path traverses WiFi, cellular, or a combination.

\textbf{Avena Messaging Network.} Built atop the overlay, this layer provides subject-based pub/sub with delay-tolerant delivery. Applications publish messages to named subjects; subscribers express interest in subject patterns with delivery requirements (priority, retention policy, ordering, geo-bounds). The messaging layer routes messages toward interested subscribers, storing at intermediate relay nodes when direct paths are unavailable and forwarding opportunistically when connectivity is established---including via mobile devices that physically transport data between disconnected segments.

The messaging network answers: \emph{how do messages reach interested subscribers given intermittent connectivity and heterogeneous costs?} It implements the core delay-tolerant pub/sub semantics that distinguish Avena from systems assuming reliable networks.

\textbf{Avena Fleet Management (avenad).} The fleet management daemon handles workload orchestration: what software runs on which devices. Workload specifications declare deployment targets, container images, environment configuration, and authorization scopes. Specifications propagate via gossip; each device reconciles local state against received specifications. Certificate management, device enrollment, and capability delegation are fleet management concerns.

Fleet management is portable: avenad uses the Messaging Network's pub/sub API and works identically on traditional networks (direct NATS) or the Avena mesh. The same binary runs on edge devices and cloud servers.

\subsection{Layered Architecture}

These three subsystems map to a four-layer stack:

\begin{Verbatim}[frame=single,fontsize=\small]
+------------------------------------------------------------------+
|  Layer 4: Fleet Management (avenad)                              |
|  - Workload specification, deployment, reconciliation            |
|  - Certificate lifecycle, capability delegation                  |
|  - Device registry, enrollment                                   |
+------------------------------------------------------------------+
|  Layer 3: State Synchronization (Gossip)                         |
|  - CRDT-based distributed state (workloads, devices, interests)  |
|  - Cost metadata propagation                                     |
|  - Shared by Messaging Network and Fleet Management              |
+------------------------------------------------------------------+
|  Layer 2: Avena Messaging Network                                |
|  - Subject-based pub/sub (NATS-compatible API)                   |
|  - Interest-driven, priority-aware routing                       |
|  - Store-and-forward for disconnected subscribers                |
|  - Geo-spatial filtering                                         |
+------------------------------------------------------------------+
|  Layer 1: Avena Overlay Network                                  |
|  - Wireguard interface per (peer, underlay) tuple                |
|  - Babel routing over all WG interfaces (unicast mode)           |
|  - LoRa broadcast mesh for low-rate control channel              |
+------------------------------------------------------------------+
|  Layer 0: Physical Links                                         |
|  - WiFi, cellular, ethernet, satellite (high-rate)               |
|  - LoRa, TV whitespace (low-rate, long-range)                    |
+------------------------------------------------------------------+
\end{Verbatim}

Layer 3 (Gossip) is the mechanism shared by Layers 2 and 4: it propagates both messaging metadata (active interests, routing costs) and fleet state (workload specifications, device information). This unification is intentional---both require the same consistency properties and benefit from the same delay-tolerant propagation.

\subsection{Workload Network Architecture}

Containerized workloads run as isolated tenants on host devices. The host's avenad daemon owns Layer 1 infrastructure; workloads interact with the network through one of two modes.

\textbf{Messaging mode} (default): The workload connects to the host avenad's local NATS-compatible endpoint. All communication occurs via pub/sub---the workload publishes messages and subscribes to subjects. The host avenad handles all overlay networking: tunnel management, routing, and message delivery. The workload has no overlay IP address and no direct network access.

\begin{Verbatim}[frame=single,fontsize=\small]
Host Device
+------------------------------------------+
|  avenad (root)                           |
|  - Owns Wireguard interfaces per underlay|
|  - Runs Babel routing                    |
|  - Exposes NATS endpoint to workloads    |
|  - Routes messages via Layer 2           |
+------------------------------------------+
       |  NATS API
       v
+------------------+  +------------------+
|  workload-A      |  |  workload-B      |
|  (Messaging)     |  |  (Messaging)     |
|  - pub/sub only  |  |  - pub/sub only  |
+------------------+  +------------------+
\end{Verbatim}

\textbf{Addressable mode}: For workloads requiring direct IP connectivity, the host avenad spawns an \texttt{avena-sidecar} container alongside the workload. The sidecar runs userspace Wireguard endpoints (one per underlay interface) with the workload's derived identity and overlay IP address. Remote nodes establish Wireguard tunnels directly to the sidecar. The host routes Wireguard UDP packets to the appropriate sidecar but cannot decrypt the traffic---tunnel keys are held only by the sidecar and remote peers.

\begin{Verbatim}[frame=single,fontsize=\small]
Host Device
+------------------------------------------+
|  avenad (root)                           |
|  - Routes Wireguard UDP to sidecars      |
|  - Announces workload IPs via Babel      |
|  - Cannot decrypt workload tunnels       |
+------------------------------------------+
       |  UDP routing
       v
+------------------+  +------------------+
|  workload-C pod  |  |  workload-D pod  |
|  +-----------+   |  |  +-----------+   |
|  | sidecar   |   |  |  | sidecar   |   |
|  | (WG keys) |   |  |  | (WG keys) |   |
|  +-----------+   |  |  +-----------+   |
|  | app       |   |  |  | app       |   |
|  +-----------+   |  |  +-----------+   |
+------------------+  +------------------+
      fd00::1122         fd00::3344
\end{Verbatim}

\textbf{Workload isolation:} In both modes, workloads are isolated from each other. Messaging-mode workloads communicate only via Layer 2 pub/sub with message-level encryption. Addressable-mode workloads have end-to-end encrypted tunnels---even workloads on the same host cannot observe each other's traffic. The host avenad can observe metadata (source/destination IPs, message timing, sizes) but not payload content.

\textbf{Overlay IP assignment:} Addressable workloads receive IPv6 addresses in the overlay's ULA range, derived from the workload's identity. The host avenad announces these addresses via Babel, advertising them as reachable through itself.

\textbf{Internet access:} Addressable workloads may require connectivity beyond the overlay (e.g., pulling container images, accessing cloud APIs). The host provides NAT for non-overlay traffic: packets destined for overlay addresses route through Wireguard tunnels via Babel, while packets destined for external addresses are forwarded to the host's internet-facing interface with source NAT applied.
